\chapter{Chapter02 — Gaussian Disagreement}
%%% generated by the blueprint pipeline from FABL.Chapter02.NoiseStability.GaussianDisagreement
\section{Overview}

Declarations of the Lean module \texttt{FABL.Chapter02.NoiseStability.GaussianDisagreement}.

\section{Declarations}

\begin{definition}[First coordinate in the correlation plane]
\lean{FABL.correlationFirstCoordinate}
\label{FABL.correlationFirstCoordinate}
\leanok
\uses{FABL.CorrelationPlane}
For a point $z=(z_1,z_2)$ in the real correlation plane, its first coordinate is $z_1$.
\end{definition}

\begin{definition}[Second coordinate in the correlation plane]
\lean{FABL.correlationSecondCoordinate}
\label{FABL.correlationSecondCoordinate}
\leanok
\uses{FABL.CorrelationPlane}
For a point $z$ in the real plane equipped with the usual Euclidean norm, its second
coordinate is the second component of its underlying pair in $\mathbb{R}^2$.
\end{definition}

\begin{theorem}[Continuity of the first coordinate]
\lean{FABL.continuous_correlationFirstCoordinate}
\label{FABL.continuous_correlationFirstCoordinate}
\leanok
\uses{FABL.correlationFirstCoordinate}
\difficulty{3}
On the real plane equipped with the $2$-norm, the map taking a point to its first real
coordinate is continuous.
\end{theorem}

\begin{theorem}[Continuity of the second coordinate]
\lean{FABL.continuous_correlationSecondCoordinate}
\label{FABL.continuous_correlationSecondCoordinate}
\leanok
\uses{FABL.correlationSecondCoordinate}
\difficulty{3}
On the real plane equipped with the $2$-norm, the map taking a point to its second real
coordinate is continuous.
\end{theorem}

\begin{definition}[Gaussian disagreement region]
\lean{FABL.gaussianDisagreementRegion}
\statementsource{boolean-ch11-gaussian-invariance}{pdf-8b2c50cf9621-p016-b001}
\label{FABL.gaussianDisagreementRegion}
\leanok
\uses{FABL.CorrelationPlane, FABL.correlationFirstCoordinate, FABL.correlationSecondCoordinate}
In the two-dimensional Euclidean plane, the Gaussian disagreement region consists of
points $z$ for which the product of the two coordinates of $z$ is strictly negative.
\end{definition}

\begin{theorem}[Openness of the Gaussian disagreement region]
\lean{FABL.isOpen_gaussianDisagreementRegion}
\label{FABL.isOpen_gaussianDisagreementRegion}
\leanok
\uses{FABL.gaussianDisagreementRegion}
\difficulty{3}
In the Euclidean plane, the set of points $(x,y)$ satisfying $xy<0$ is open.
\end{theorem}

\begin{theorem}[Boundary of the disagreement region]
\lean{FABL.frontier_gaussianDisagreementRegion_subset}
\label{FABL.frontier_gaussianDisagreementRegion_subset}
\leanok
\uses{FABL.CorrelationPlane, FABL.correlationFirstCoordinate, FABL.correlationSecondCoordinate, FABL.gaussianDisagreementRegion}
\difficulty{3}
In the real coordinate plane, let the disagreement region consist of points whose two
coordinates have product less than zero. Every point on its frontier has at least one
coordinate equal to zero.
\end{theorem}

\begin{theorem}[First coordinate has standard Gaussian law]
\lean{FABL.correlatedGaussianMeasure_map_firstCoordinate}
\label{FABL.correlatedGaussianMeasure_map_firstCoordinate}
\leanok
\uses{FABL.CorrelationPlane, FABL.continuous_correlationFirstCoordinate, FABL.continuous_gaussianCorrelationMap, FABL.correlatedGaussianMeasure, FABL.correlationFirstCoordinate, FABL.gaussianCorrelationMap}
\difficulty{4}
For every $\rho\in\mathbb R$, if $(X,Y)$ has the probability distribution obtained from
independent standard Gaussian variables $U,V$ by setting $X=U$ and $Y=\rho
U+\sqrt{1-\rho^2}\,V$, then the marginal distribution of $X$ is the standard Gaussian
distribution on $\mathbb R$.
\end{theorem}

\begin{theorem}[Second marginal of a correlated Gaussian]
\lean{FABL.correlatedGaussianMeasure_map_secondCoordinate}
\label{FABL.correlatedGaussianMeasure_map_secondCoordinate}
\leanok
\uses{FABL.CorrelationPlane, FABL.charFun_correlatedGaussianMeasure, FABL.continuous_correlationSecondCoordinate, FABL.correlatedGaussianMeasure, FABL.correlationQuadraticForm, FABL.correlationSecondCoordinate}
\difficulty{4}
Let $\rho\in[-1,1]$, and let $(X,Y)$ have the law obtained by taking independent
standard normal random variables $U,V$ and setting $X=U$ and $Y=\rho
U+\sqrt{1-\rho^2}\,V$. Then the second coordinate $Y$ has the standard normal
distribution.
\end{theorem}

\begin{definition}[Standard Gaussian density on the plane]
\lean{FABL.standardGaussianPlaneDensity}
\label{FABL.standardGaussianPlaneDensity}
\leanok
For $u,v\in\mathbb{R}$, define the standard Gaussian density on the plane by
$\rho(u,v)=\frac{1}{2\pi}\exp\!\left(-\frac{u^2+v^2}{2}\right)$, viewed as a nonnegative
extended-real-valued function.
\end{definition}

\begin{definition}[Standard Gaussian radial density]
\lean{FABL.standardGaussianPolarRadialDensity}
\label{FABL.standardGaussianPolarRadialDensity}
\leanok
For $r\in\mathbb{R}$, define the nonnegative extended real-valued function $\rho$ by
$\rho(r)=\max\{r,0\}(2\pi)^{-1}e^{-r^2/2}$.
\end{definition}

\begin{theorem}[Polar product identity for Gaussian densities]
\lean{FABL.gaussianPDF_mul_polar}
\label{FABL.gaussianPDF_mul_polar}
\leanok
\difficulty{4}
For all real numbers $r$ and $\theta$, the product of the standard Gaussian densities at
$r\cos\theta$ and $r\sin\theta$ equals the product of the standard Gaussian densities at
$r$ and at $0$.
\end{theorem}

\begin{theorem}[Gaussian plane density in polar coordinates]
\lean{FABL.standardGaussianPlaneDensity_polar}
\label{FABL.standardGaussianPlaneDensity_polar}
\leanok
\uses{FABL.gaussianPDF_mul_polar, FABL.standardGaussianPlaneDensity, FABL.standardGaussianPolarRadialDensity}
\difficulty{3}
For every pair $p=(r,\theta)\in\mathbb{R}\times\mathbb{R}$, multiplying the standard
Gaussian plane density at the point with polar coordinates $(r,\theta)$ by $r$ gives the
standard Gaussian polar radial density at $r$.
\end{theorem}

\begin{theorem}[Measurability of the Gaussian plane density]
\lean{FABL.measurable_standardGaussianPlaneDensity}
\label{FABL.measurable_standardGaussianPlaneDensity}
\leanok
\uses{FABL.standardGaussianPlaneDensity}
\difficulty{2}
The function on $\mathbb{R}^2$ given by $(u,v)\mapsto \varphi(u)\varphi(v)$ is
measurable, where $\varphi$ is the standard Gaussian density on $\mathbb{R}$.
\end{theorem}

\begin{theorem}[Measurability of the radial density]
\lean{FABL.measurable_standardGaussianPolarRadialDensity}
\label{FABL.measurable_standardGaussianPolarRadialDensity}
\leanok
\uses{FABL.standardGaussianPolarRadialDensity}
\difficulty{2}
The function on $\mathbb{R}$ given by $r\mapsto \max(r,0)\,\varphi(r)\varphi(0)$ is
measurable, where $\varphi$ is the standard Gaussian density on $\mathbb{R}$.
\end{theorem}

\begin{theorem}[Product Gaussian measure has the plane density]
\lean{FABL.standardGaussianPlane_eq_withDensity}
\label{FABL.standardGaussianPlane_eq_withDensity}
\leanok
\uses{FABL.standardGaussianPlaneDensity}
\difficulty{3}
The product of two standard Gaussian measures on $\mathbb{R}$ equals the measure on
$\mathbb{R}^2$ obtained by weighting Lebesgue measure by the standard Gaussian plane
density.
\end{theorem}

\begin{theorem}[Angular events under the planar Gaussian measure]
\lean{FABL.standardGaussianPlane_measure_angularEvent}
\label{FABL.standardGaussianPlane_measure_angularEvent}
\leanok
\uses{FABL.measurable_standardGaussianPolarRadialDensity, FABL.standardGaussianPlaneDensity, FABL.standardGaussianPlaneDensity_polar, FABL.standardGaussianPlane_eq_withDensity, FABL.standardGaussianPolarRadialDensity}
\difficulty{5}
Let $s\subseteq\mathbb{R}^2$ and let $A\subseteq\mathbb{R}$ be measurable, with
$A\subseteq(-\pi,\pi)$. Suppose that for every polar coordinate pair $p=(r,\theta)$ with
$r>0$ and $-\pi<\theta<\pi$, the point with polar coordinates $(r,\theta)$ lies in $s$
exactly when $\theta\in A$. Then the product of two standard Gaussian measures assigns
to $s$ the measure $\left(\int_0^\infty r\,\varphi(r)\varphi(0)\,dr\right)|A|$, where
$\varphi$ is the standard Gaussian density and $|A|$ is the Lebesgue measure of $A$.
\end{theorem}

\begin{definition}[Standard Gaussian polar radial mass]
\lean{FABL.standardGaussianPolarRadialMass}
\label{FABL.standardGaussianPolarRadialMass}
\leanok
\uses{FABL.standardGaussianPolarRadialDensity}
For each natural number $n$, define the standard Gaussian polar radial mass to be the
integral of the standard Gaussian polar radial density over positive radii: $\int_{r>0}
\text{standard Gaussian polar radial density in dimension }n\,dr$, interpreted as a
nonnegative extended real number.
\end{definition}

\begin{theorem}[Total radial mass for the planar Gaussian]
\lean{FABL.standardGaussianPolarRadialMass_eq}
\label{FABL.standardGaussianPolarRadialMass_eq}
\leanok
\uses{FABL.standardGaussianPlane_measure_angularEvent, FABL.standardGaussianPolarRadialMass}
\difficulty{3}
The total radial mass of the standard Gaussian in polar coordinates is $(2\pi)^{-1}$.
\end{theorem}

\begin{definition}[Gaussian source disagreement region]
\lean{FABL.gaussianSourceDisagreement}
\statementsource{boolean-ch11-gaussian-invariance}{pdf-8b2c50cf9621-p016-b001}
\label{FABL.gaussianSourceDisagreement}
\leanok
For each real number $\rho$, define the Gaussian source disagreement region to be the
set of pairs $(u,v)\in\mathbb R^2$ such that $u\bigl(\rho u+\sqrt{1-\rho^2}\,v\bigr)<0$.
\end{definition}

\begin{theorem}[Openness of the Gaussian source disagreement set]
\lean{FABL.isOpen_gaussianSourceDisagreement}
\label{FABL.isOpen_gaussianSourceDisagreement}
\leanok
\uses{FABL.gaussianSourceDisagreement}
\difficulty{3}
For each real number $\rho$, the set of points $(u,v)\in\mathbb R^2$ satisfying
$u\bigl(\rho u+\sqrt{1-\rho^2}\,v\bigr)<0$ is open.
\end{theorem}

\begin{theorem}[Correlated gaussian measure disagreement eq source]
\lean{FABL.correlatedGaussianMeasure_disagreement_eq_source}
\statementsource{boolean-ch11-gaussian-invariance}{pdf-8b2c50cf9621-p016-b001}
\label{FABL.correlatedGaussianMeasure_disagreement_eq_source}
\leanok
\uses{FABL.CorrelationPlane, FABL.continuous_gaussianCorrelationMap, FABL.correlatedGaussianMeasure, FABL.gaussianCorrelationMap, FABL.gaussianDisagreementRegion, FABL.gaussianSourceDisagreement, FABL.isOpen_gaussianDisagreementRegion}
\difficulty{3}
% dataset-pipeline: stub -- prose pending build_informal_statements.py
Statement pending regeneration.
\end{theorem}

\begin{definition}[Angular disagreement set]
\lean{FABL.angularDisagreementSet}
\label{FABL.angularDisagreementSet}
\leanok
For $\alpha\in\mathbb{R}$, the angular disagreement set consists of the angles
$\theta\in(-\pi,\pi)$ for which $\cos\theta\,\cos(\theta-\alpha)<0$.
\end{definition}

\begin{theorem}[Openness of the angular disagreement set]
\lean{FABL.isOpen_angularDisagreementSet}
\label{FABL.isOpen_angularDisagreementSet}
\leanok
\uses{FABL.angularDisagreementSet}
\difficulty{3}
For each real number $\alpha$, the set of angles $\theta\in(-\pi,\pi)$ for which
$\cos\theta\,\cos(\theta-\alpha)<0$ is open in $\mathbb R$.
\end{theorem}

\begin{theorem}[Positive cosine on a principal interval]
\lean{FABL.cos_pos_iff_of_neg_pi_lt_of_lt_pi}
\label{FABL.cos_pos_iff_of_neg_pi_lt_of_lt_pi}
\leanok
\difficulty{4}
For every real number $x$ with $-\pi<x<\pi$, $\cos x>0$ if and only if $-\pi/2<x<\pi/2$.
\end{theorem}

\begin{theorem}[Negative cosine on a principal interval]
\lean{FABL.cos_neg_iff_of_neg_pi_lt_of_lt_pi}
\label{FABL.cos_neg_iff_of_neg_pi_lt_of_lt_pi}
\leanok
\difficulty{4}
For every real number $x$ with $-\pi<x<\pi$, $\cos x<0$ if and only if $x< -\pi/2$ or
$\pi/2<x$.
\end{theorem}

\begin{theorem}[Positive cosine on an extended interval]
\lean{FABL.cos_pos_iff_of_neg_two_pi_lt_of_lt_pi}
\label{FABL.cos_pos_iff_of_neg_two_pi_lt_of_lt_pi}
\leanok
\uses{FABL.cos_pos_iff_of_neg_pi_lt_of_lt_pi}
\difficulty{5}
For every real number $x$ with $-2\pi<x<\pi$, $\cos x>0$ if and only if either
$x<-3\pi/2$ or $-\pi/2<x<\pi/2$.
\end{theorem}

\begin{theorem}[Negative cosine on an extended interval]
\lean{FABL.cos_neg_iff_of_neg_two_pi_lt_of_lt_pi}
\label{FABL.cos_neg_iff_of_neg_two_pi_lt_of_lt_pi}
\leanok
\uses{FABL.cos_neg_iff_of_neg_pi_lt_of_lt_pi}
\difficulty{5}
For every real number $x$ with $-2\pi<x<\pi$, $\cos x<0$ if and only if either
$-3\pi/2<x<-\pi/2$ or $\pi/2<x$.
\end{theorem}

\begin{theorem}[Angular disagreement intervals]
\lean{FABL.angularDisagreementSet_eq_of_le_half_pi}
\label{FABL.angularDisagreementSet_eq_of_le_half_pi}
\leanok
\uses{FABL.angularDisagreementSet, FABL.cos_neg_iff_of_neg_pi_lt_of_lt_pi, FABL.cos_neg_iff_of_neg_two_pi_lt_of_lt_pi, FABL.cos_pos_iff_of_neg_pi_lt_of_lt_pi, FABL.cos_pos_iff_of_neg_two_pi_lt_of_lt_pi}
\difficulty{5}
For $0\leq\alpha\leq\pi/2$, the set of angles $\theta\in(-\pi,\pi)$ for which
$\cos\theta\cos(\theta-\alpha)<0$ is
$(-\pi/2,\alpha-\pi/2)\cup(\pi/2,\alpha+\pi/2)$.
\end{theorem}

\begin{theorem}[Angular disagreement measure for small angles]
\lean{FABL.volume_angularDisagreementSet_of_le_half_pi}
\label{FABL.volume_angularDisagreementSet_of_le_half_pi}
\leanok
\uses{FABL.angularDisagreementSet, FABL.angularDisagreementSet_eq_of_le_half_pi}
\difficulty{4}
For $0\leq\alpha\leq\pi/2$, consider the angles $\theta\in(-\pi,\pi)$ for which
$\cos\theta$ and $\cos(\theta-\alpha)$ have opposite signs. Their Lebesgue measure is
$2\alpha$.
\end{theorem}

\begin{theorem}[Angular disagreement intervals]
\lean{FABL.angularDisagreementSet_eq_of_half_pi_le}
\label{FABL.angularDisagreementSet_eq_of_half_pi_le}
\leanok
\uses{FABL.angularDisagreementSet, FABL.cos_neg_iff_of_neg_pi_lt_of_lt_pi, FABL.cos_neg_iff_of_neg_two_pi_lt_of_lt_pi, FABL.cos_pos_iff_of_neg_pi_lt_of_lt_pi, FABL.cos_pos_iff_of_neg_two_pi_lt_of_lt_pi}
\difficulty{5}
For $\pi/2\leq\alpha\leq\pi$, the set of angles $\theta\in(-\pi,\pi)$ for which
$\cos\theta\cos(\theta-\alpha)<0$ is
$(-\pi,\alpha-3\pi/2)\cup(-\pi/2,\alpha-\pi/2)\cup(\pi/2,\pi)$.
\end{theorem}

\begin{theorem}[Angular disagreement measure for large angles]
\lean{FABL.volume_angularDisagreementSet_of_half_pi_le}
\label{FABL.volume_angularDisagreementSet_of_half_pi_le}
\leanok
\uses{FABL.angularDisagreementSet, FABL.angularDisagreementSet_eq_of_half_pi_le}
\difficulty{4}
For $\pi/2\leq\alpha\leq\pi$, consider the angles $\theta\in(-\pi,\pi)$ for which
$\cos\theta$ and $\cos(\theta-\alpha)$ have opposite signs. Their Lebesgue measure is
$2\alpha$.
\end{theorem}

\begin{theorem}[Angular disagreement measure]
\lean{FABL.volume_angularDisagreementSet}
\label{FABL.volume_angularDisagreementSet}
\leanok
\uses{FABL.angularDisagreementSet, FABL.volume_angularDisagreementSet_of_half_pi_le, FABL.volume_angularDisagreementSet_of_le_half_pi}
\difficulty{4}
For $0\leq\alpha\leq\pi$, consider the angles $\theta\in(-\pi,\pi)$ for which
$\cos\theta$ and $\cos(\theta-\alpha)$ have opposite signs. Their Lebesgue measure is
$2\alpha$.
\end{theorem}

\begin{theorem}[Polar characterization of Gaussian disagreement]
\lean{FABL.gaussianSourceDisagreement_polar}
\label{FABL.gaussianSourceDisagreement_polar}
\leanok
\uses{FABL.angularDisagreementSet, FABL.gaussianSourceDisagreement}
\difficulty{5}
Let $\rho\in[-1,1]$ and let $p=(r,\theta)$ lie in the target of the polar-coordinate
map. The point obtained by applying the inverse polar-coordinate map to $p$ belongs to
the Gaussian source disagreement set at correlation $\rho$ if and only if
$\theta\in(-\pi,\pi)$ and $\cos\theta\,\cos(\theta-\arccos\rho)<0$.
\end{theorem}

\begin{theorem}[Correlated gaussian measure disagreement]
\lean{FABL.correlatedGaussianMeasure_disagreement}
\statementsource{Analysis-of-Boolean-Functions-by-Ryan-ODonnell}{
  pdf-ffbe9c7b3306-p120-b003,
  pdf-ffbe9c7b3306-p120-b004
}
\label{FABL.correlatedGaussianMeasure_disagreement}
\leanok
\uses{FABL.CorrelationPlane, FABL.angularDisagreementSet, FABL.correlatedGaussianMeasure, FABL.correlatedGaussianMeasure_disagreement_eq_source, FABL.gaussianDisagreementRegion, FABL.gaussianSourceDisagreement, FABL.gaussianSourceDisagreement_polar, FABL.isOpen_angularDisagreementSet, FABL.isOpen_gaussianSourceDisagreement, FABL.standardGaussianPlane_measure_angularEvent, FABL.standardGaussianPolarRadialMass, FABL.standardGaussianPolarRadialMass_eq, FABL.volume_angularDisagreementSet}
\difficulty{4}
% dataset-pipeline: stub -- prose pending build_informal_statements.py
Statement pending regeneration.
\end{theorem}

\begin{theorem}[Null coordinate axes under correlated Gaussian measure]
\lean{FABL.correlatedGaussianMeasure_coordinateAxes_null}
\label{FABL.correlatedGaussianMeasure_coordinateAxes_null}
\leanok
\uses{FABL.CorrelationPlane, FABL.continuous_correlationFirstCoordinate, FABL.continuous_correlationSecondCoordinate, FABL.correlatedGaussianMeasure, FABL.correlatedGaussianMeasure_map_firstCoordinate, FABL.correlatedGaussianMeasure_map_secondCoordinate, FABL.correlationFirstCoordinate, FABL.correlationSecondCoordinate}
\difficulty{4}
Let $-1\le\rho\le1$, and let $(X,Y)$ have the distribution obtained by mapping two
independent standard normal variables $(U,V)$ to $(U,\rho U+\sqrt{1-\rho^2}V)$. The
probability that either coordinate of $(X,Y)$ is zero is zero.
\end{theorem}

\begin{theorem}[Nullity of the disagreement boundary]
\lean{FABL.correlatedGaussianMeasure_frontier_disagreement_null}
\label{FABL.correlatedGaussianMeasure_frontier_disagreement_null}
\leanok
\uses{FABL.CorrelationPlane, FABL.correlatedGaussianMeasure, FABL.correlatedGaussianMeasure_coordinateAxes_null, FABL.frontier_gaussianDisagreementRegion_subset, FABL.gaussianDisagreementRegion}
\difficulty{1}
Let $-1\leq\rho\leq1$, and let $(X,Y)$ have the probability distribution obtained from
independent standard Gaussian variables $U,V$ by setting $X=U$ and $Y=\rho
U+\sqrt{1-\rho^2}\,V$. The boundary of the region $\{(x,y):xy<0\}$ has probability zero
under this distribution.
\end{theorem}

\begin{theorem}[Normalized correlated sign sum measure disagreement tendsto]
\lean{FABL.normalizedCorrelatedSignSumMeasure_disagreement_tendsto}
\statementsource{Analysis-of-Boolean-Functions-by-Ryan-ODonnell}{
  pdf-ffbe9c7b3306-p120-b003,
  pdf-ffbe9c7b3306-p054-b005
}
\statementsource{mossel-neeman-gaussian-noise-stability}{
  pdf-b53c1a5aa27d-p007-b003,
  pdf-b53c1a5aa27d-p008-b001
}
\label{FABL.normalizedCorrelatedSignSumMeasure_disagreement_tendsto}
\leanok
\uses{FABL.correlatedGaussianMeasure, FABL.correlatedGaussianMeasure_frontier_disagreement_null, FABL.gaussianDisagreementRegion, FABL.normalizedCorrelatedSignSumMeasure, FABL.normalizedCorrelatedSignSumMeasure_tendsto_correlatedGaussian}
\difficulty{1}
% dataset-pipeline: stub -- prose pending build_informal_statements.py
Statement pending regeneration.
\end{theorem}

\begin{theorem}[First coordinate of the normalized pair sum]
\lean{FABL.correlationFirstCoordinate_normalizedCorrelatedPairSum}
\label{FABL.correlationFirstCoordinate_normalizedCorrelatedPairSum}
\leanok
\uses{FABL.correlatedSignVector, FABL.correlationFirstCoordinate, FABL.normalizedCorrelatedPairSum, FABL.normalizedCorrelatedSignVectorSum, FABL.pairCoordinatesEquiv, FABL.signValue}
\difficulty{3}
For every natural number $n$ and every pair of sign vectors $x,y$ of length $n$, the
first coordinate of their normalized correlated sum is $(\sqrt n)^{-1}\sum_{i=1}^n x_i$,
with signs interpreted as real numbers.
\end{theorem}

\begin{theorem}[Second coordinate of a normalized pair sum]
\lean{FABL.correlationSecondCoordinate_normalizedCorrelatedPairSum}
\label{FABL.correlationSecondCoordinate_normalizedCorrelatedPairSum}
\leanok
\uses{FABL.correlatedSignVector, FABL.correlationSecondCoordinate, FABL.normalizedCorrelatedPairSum, FABL.normalizedCorrelatedSignVectorSum, FABL.pairCoordinatesEquiv, FABL.signValue}
\difficulty{3}
For every natural number $n$ and every pair $x,y\in\{ -1,1\}^n$ of sign vectors, the
second coordinate of their normalized correlated sum is $(\sqrt n)^{-1}\sum_{i=1}^n
y_i$.
\end{theorem}

\begin{theorem}[Opposite threshold signs]
\lean{FABL.thresholdSign_ne_iff_mul_neg}
\label{FABL.thresholdSign_ne_iff_mul_neg}
\leanok
\uses{FABL.thresholdSign, FABL.thresholdSign_of_neg, FABL.thresholdSign_of_nonneg}
\difficulty{5}
For nonzero real numbers $a$ and $b$, their threshold signs differ if and only if
$ab<0$.
\end{theorem}

\begin{theorem}[Majority disagreement and the disagreement region]
\lean{FABL.normalizedCorrelatedPairSum_mem_disagreement_iff_majority_ne}
\label{FABL.normalizedCorrelatedPairSum_mem_disagreement_iff_majority_ne}
\leanok
\uses{FABL.correlationFirstCoordinate, FABL.correlationFirstCoordinate_normalizedCorrelatedPairSum, FABL.correlationSecondCoordinate, FABL.correlationSecondCoordinate_normalizedCorrelatedPairSum, FABL.gaussianDisagreementRegion, FABL.majority, FABL.normalizedCorrelatedPairSum, FABL.signValue, FABL.sum_signValue_ne_zero_of_odd, FABL.thresholdSign, FABL.thresholdSign_ne_iff_mul_neg}
\difficulty{4}
Let $n$ be odd, and let $x,y\in\{-1,1\}^n$. Form the point in the plane whose
coordinates are the respective sums of the signs in $x$ and $y$, each divided by $\sqrt
n$. This point has coordinates with negative product if and only if the majorities of
$x$ and $y$ differ, where the majority sign is positive when the sum of signs is
nonnegative and negative otherwise.
\end{theorem}
